\begin{pastexam}[label={px:ust-112-ee-ds-10}]{臺聯大電機類 112 學年度 第 10 題}
（單選；3 分）

Given the following code, what will happen during the execution?
\begin{lstlisting}[basicstyle=\ttfamily\small, frame=single, columns=fullflexible,
    keepspaces=true, xleftmargin=6em, xrightmargin=6em]
1  void do_job()
2  {
3      do_job();
4  }
5  int main()
6  {
7      do_job;
8      return 0;
9  }
\end{lstlisting}
(A)~The code cannot be compiled\\
(B)~The code will be executed successfully and ``0'' will be printed out\\
(C)~The code will be executed successfully and random output will be exported\\
(D)~The code will be executed successfully, but an overflow will occur
\end{pastexam}

\begin{pastsoln}
答案：照出題意圖是 (D)；照字面，四個選項都不完全對。

先說字面。 第 7 行寫的是 \texttt{do\_job;}，沒有括號。 在 C 裡，函式名稱單獨出現是一個運算式
（值是函式的位址），\texttt{do\_job;} 是一個「算完就丟掉」的運算式敘述，合法，編譯器最多給
「statement with no effect」的警告。 它\textbf{不會呼叫} \texttt{do\_job}。 所以照字面：
\begin{itemize}
\item (A)：錯，可以編譯。
\item (B)：錯。 \texttt{return 0} 是把 $0$ 當成程式的結束狀態交給作業系統，不是印出來；程式
沒有任何輸出敘述。
\item (C)：錯，沒有任何輸出。
\item (D)：錯，\texttt{do\_job} 根本沒被呼叫，不會溢位。
\end{itemize}
照字面，程式正常結束、什麼也不印，沒有對應的選項。

題目的考點顯然是第 1--4 行：\texttt{do\_job} 無條件呼叫自己，沒有終止條件。 把第 7 行讀成
\texttt{do\_job();}（漏打括號），執行時每呼叫一次就在呼叫堆疊上多一個活動紀錄（返回位址等），
永遠不會返回，堆疊空間很快用完，發生 stack overflow，程式被作業系統終止（Linux 上通常是
segmentation fault）（資結書第~6.4~節）。 這和 (D) 說的「overflow will occur」一致，所以這樣讀時
答案是 (D)。 (D) 前半句「executed successfully」的意思只能理解成「可以編譯、開始執行」，程式
實際上是異常結束的。

另外，有些編譯器在最佳化時會把這種尾端的自我呼叫改成迴圈，那時程式會變成無窮迴圈而不是
溢位；題目的意思是一般（不最佳化）的情形。
\end{pastsoln}
